% ============================================================
%  Section 2.3.3 – Description textuelle des cas d'utilisation
%  Projet : BOA AI Hotline Automation
% ============================================================

% Packages requis (à ajouter au préambule s'ils ne le sont pas déjà)
% \usepackage[table,xcdraw]{xcolor}
% \usepackage{tabularx}
% \usepackage{tcolorbox}
% \usepackage{fontawesome5}
% \usepackage{float}          % REQUIS pour forcer la position avec [H]

% ── Couleurs BOA ────────────────────────────────────────────
\definecolor{boaGreen}{HTML}{2E7D32}       % vert foncé entête
\definecolor{boaLightGreen}{HTML}{C8E6C9}  % vert clair lignes
\definecolor{boaGray}{HTML}{F5F5F5}        % gris très clair lignes
\definecolor{guidelineBg}{HTML}{F1F8E9}    % fond bloc guideline
\definecolor{guidelineBorder}{HTML}{558B2F}% bordure bloc guideline

% ─────────────────────────────────────────────────────────────
\subsection{Description textuelle des cas d'utilisation}
% ─────────────────────────────────────────────────────────────

% ── Bloc Guideline ──────────────────────────────────────────
\begin{tcolorbox}[
  colback   = guidelineBg,
  colframe  = guidelineBorder,
  left      = 6pt,
  right     = 6pt,
  top       = 4pt,
  bottom    = 4pt,
  arc       = 2pt,
  boxrule   = 0.8pt
]
  \textbf{\faBookmark\; Guideline :}
  \textit{Pour chaque cas d'utilisation principal, fournissez une description textuelle structurée selon le tableau ci-dessous.}
\end{tcolorbox}

\vspace{10pt}

% ── CU 1 : Poser une question (FAQ) ──────────────────────────
\begin{table}[H]
  \centering
  \renewcommand{\arraystretch}{1.5}
  \setlength{\tabcolsep}{10pt}
  \begin{tabularx}{\textwidth}{|>{\bfseries\color{white}\columncolor{boaGreen}}l|X|}
    \hline
    \multicolumn{2}{|>{\bfseries\color{white}\columncolor{boaGreen}}c|}{Cas d'utilisation : Poser une question et obtenir une réponse FAQ (CU-01)} \\ \hline
    Acteurs & Client BOA / Utilisateur \\ \hline
    Pré-conditions & 
    Le client a démarré l'interface de discussion (ChatBot).\\ \hline
    Scénario nominal & 
    1. Le client pose une question ou saisit une requête dans le ChatBot.\newline
    2. Le ChatBot authentifie la session de l'utilisateur.\newline
    3. Le système analyse l'intention de la question via le module IA de classification.\newline
    4. Le système récupère la réponse adéquate dans la FAQ.\newline
    5. Le système affiche la réponse sur l'interface du ChatBot.\newline
    6. Le client consulte la réponse. \\ \hline
    Scénarios alternatifs & 
    \textbf{A3. Intention non reconnue :}\newline
    Si le module IA ne comprend pas la question, le ChatBot demande au client de reformuler sa requête.\newline
    \textbf{A4. Réponse non trouvée dans la FAQ :}\newline
    Si la base de connaissances ne contient pas la réponse, le ChatBot propose au client de contacter le support ou de formuler une demande d'assistance. \\ \hline
    Post-conditions & Le client a reçu la réponse à sa question et la discussion se poursuit. \\ \hline
  \end{tabularx}
  \caption{Description du cas d'utilisation - Poser une question (FAQ)}
  \label{tab:cu_faq}
\end{table}

\vspace{15pt}

% ── CU 2 : Faire un virement ──────────────────────────────────
\begin{table}[H]
  \centering
  \renewcommand{\arraystretch}{1.5}
  \setlength{\tabcolsep}{10pt}
  \begin{tabularx}{\textwidth}{|>{\bfseries\color{white}\columncolor{boaGreen}}l|X|}
    \hline
    \multicolumn{2}{|>{\bfseries\color{white}\columncolor{boaGreen}}c|}{Cas d'utilisation : Déclencher une tâche RPA - Faire un virement (CU-02)} \\ \hline
    Acteurs & Client BOA / Utilisateur (Initiateur), Robot ROA / UiPath (Exécuteur) \\ \hline
    Pré-conditions & 
    - Le client est connecté au ChatBot.\newline
    - Le client s'est authentifié sur le Portail BOA. \\ \hline
    Scénario nominal & 
    1. Le client formule une demande de virement auprès du ChatBot.\newline
    2. Le ChatBot sollicite la confirmation de l'authentification Portail BOA.\newline
    3. Le client saisit les informations du virement (bénéficiaire, montant).\newline
    4. Le système envoie un code de confirmation (OTP) au client.\newline
    5. Le client saisit le code de confirmation valide.\newline
    6. Le backend (Spring Boot) valide le virement et transmet la commande de déclenchement au Robot ROA (UiPath).\newline
    7. Le Robot ROA effectue le virement sur le système bancaire central.\newline
    8. Le Robot ROA retourne le statut de l'opération (succès) au backend.\newline
    9. Le ChatBot notifie le client du succès de l'opération. \\ \hline
    Scénarios alternatifs & 
    \textbf{A5. Code de confirmation invalide :}\newline
    Si le code saisi est incorrect ou expiré, le système refuse la transaction et le ChatBot demande de générer un nouveau code.\newline
    \textbf{A7. Solde insuffisant / Limite de virement atteinte :}\newline
    Le système de décision ou le Robot ROA constate un problème d'éligibilité. L'opération est annulée et le client est notifié de l'échec de la transaction.\\ \hline
    Post-conditions & Le virement est exécuté de manière sécurisée et les soldes de comptes sont mis à jour. \\ \hline
  \end{tabularx}
  \caption{Description du cas d'utilisation - Faire un virement}
  \label{tab:cu_virement}
\end{table}

\vspace{15pt}

% ── CU 3 : Débloquer Carte ───────────────────────────────────
\begin{table}[H]
  \centering
  \renewcommand{\arraystretch}{1.5}
  \setlength{\tabcolsep}{10pt}
  \begin{tabularx}{\textwidth}{|>{\bfseries\color{white}\columncolor{boaGreen}}l|X|}
    \hline
    \multicolumn{2}{|>{\bfseries\color{white}\columncolor{boaGreen}}c|}{Cas d'utilisation : Déclencher une tâche RPA - Débloquer Carte (CU-03)} \\ \hline
    Acteurs & Client BOA / Utilisateur (Initiateur), Robot ROA / UiPath (Exécuteur) \\ \hline
    Pré-conditions & 
    - Le client est connecté au ChatBot.\newline
    - Le client s'est authentifié sur le Portail BOA.\newline
    - Le client possède une carte bancaire actuellement bloquée ou restreinte. \\ \hline
    Scénario nominal & 
    1. Le client demande à débloquer sa carte bancaire via le ChatBot.\newline
    2. Le ChatBot demande confirmation de son identité via le Portail BOA.\newline
    3. Le client sélectionne la carte à débloquer.\newline
    4. Le système envoie un code de confirmation (OTP) au client.\newline
    5. Le client saisit le code de confirmation valide.\newline
    6. Le backend transmet l'ordre de déblocage au Robot ROA (UiPath).\newline
    7. Le Robot ROA accède au système monétique central et effectue le déblocage.\newline
    8. Le Robot ROA confirme le succès du déblocage au backend.\newline
    9. Le ChatBot informe le client que sa carte est à nouveau active. \\ \hline
    Scénarios alternatifs & 
    \textbf{A5. Code de confirmation (OTP) incorrect :}\newline
    La carte reste bloquée et le client doit redemander un code.\newline
    \textbf{A7. Carte non éligible au déblocage autonome (ex: carte expirée, volée) :}\newline
    Le système ou le Robot rejette la demande, et le ChatBot invite le client à contacter un conseiller physique. \\ \hline
    Post-conditions & La carte bancaire du client est débloquée avec succès et prête à l'emploi. \\ \hline
  \end{tabularx}
  \caption{Description du cas d'utilisation - Débloquer une carte bancaire}
  \label{tab:cu_debloquer_carte}
\end{table}

\vspace{15pt}

% ── CU 4 : Activer Dotation E-Commerce ───────────────────────
\begin{table}[H]
  \centering
  \renewcommand{\arraystretch}{1.5}
  \setlength{\tabcolsep}{10pt}
  \begin{tabularx}{\textwidth}{|>{\bfseries\color{white}\columncolor{boaGreen}}l|X|}
    \hline
    \multicolumn{2}{|>{\bfseries\color{white}\columncolor{boaGreen}}c|}{Cas d'utilisation : Déclencher une tâche RPA - Activer Dotation E-Commerce (CU-04)} \\ \hline
    Acteurs & Client BOA / Utilisateur (Initiateur), Robot ROA / UiPath (Exécuteur) \\ \hline
    Pré-conditions & 
    - Le client est connecté au ChatBot.\newline
    - Le client s'est authentifié sur le Portail BOA. \\ \hline
    Scénario nominal & 
    1. Le client demande à activer sa dotation e-commerce.\newline
    2. Le ChatBot demande confirmation via le Portail BOA.\newline
    3. Le client définit ou valide le montant de la dotation souhaitée.\newline
    4. Le système transmet un code de confirmation (OTP) par SMS/email.\newline
    5. Le client saisit le code de confirmation valide.\newline
    6. Le backend Spring Boot valide la demande et réveille le Robot ROA.\newline
    7. Le Robot ROA applique l'activation de la dotation sur la carte concernée dans le système d'autorisation.\newline
    8. Le Robot ROA renvoie le statut (succès) au backend.\newline
    9. Le ChatBot confirme au client l'activation immédiate de sa dotation. \\ \hline
    Scénarios alternatifs & 
    \textbf{A3. Plafond autorisé dépassé :}\newline
    Si le client demande un montant supérieur à sa limite autorisée, le système de décision rejette la demande et notifie le client du maximum possible. \\ \hline
    Post-conditions & La dotation E-Commerce est activée sur la carte bancaire du client. \\ \hline
  \end{tabularx}
  \caption{Description du cas d'utilisation - Activer la dotation E-Commerce}
  \label{tab:cu_dotation_ecommerce}
\end{table}

\vspace{15pt}

% ── CU 5 : Activer Dotation Touristique ─────────────────────
\begin{table}[H]
  \centering
  \renewcommand{\arraystretch}{1.5}
  \setlength{\tabcolsep}{10pt}
  \begin{tabularx}{\textwidth}{|>{\bfseries\color{white}\columncolor{boaGreen}}l|X|}
    \hline
    \multicolumn{2}{|>{\bfseries\color{white}\columncolor{boaGreen}}c|}{Cas d'utilisation : Déclencher une tâche RPA - Activer Dotation Touristique (CU-05)} \\ \hline
    Acteurs & Client BOA / Utilisateur (Initiateur), Robot ROA / UiPath (Exécuteur) \\ \hline
    Pré-conditions & 
    - Le client est connecté au ChatBot.\newline
    - Le client s'est authentifié sur le Portail BOA. \\ \hline
    Scénario nominal & 
    1. Le client demande l'activation de sa dotation touristique internationale.\newline
    2. Le ChatBot initie l'authentification sécurisée.\newline
    3. Le client sélectionne sa carte et indique le montant requis (dans la limite légale annuelle).\newline
    4. Le système envoie un code de confirmation OTP.\newline
    5. Le client saisit l'OTP valide.\newline
    6. Le backend transmet la tâche d'activation de dotation internationale au Robot ROA.\newline
    7. Le Robot ROA configure la dotation touristique sur la carte du client dans le système d'autorisation.\newline
    8. Le Robot ROA retourne le succès de l'activation.\newline
    9. Le ChatBot confirme au client que sa dotation est active pour ses transactions à l'étranger. \\ \hline
    Scénarios alternatifs & 
    \textbf{A3. Droits de dotation épuisés ou quota annuel atteint :}\newline
    Le système ou le Robot identifie que le client a déjà consommé son droit annuel. L'opération est refusée et le client est averti par le ChatBot. \\ \hline
    Post-conditions & La dotation touristique est allouée et activée sur le profil monétique du client. \\ \hline
  \end{tabularx}
  \caption{Description du cas d'utilisation - Activer la dotation touristique}
  \label{tab:cu_dotation_touristique}
\end{table}

\vspace{15pt}

% ── CU 6 : Administration du système ─────────────────────────
\begin{table}[H]
  \centering
  \renewcommand{\arraystretch}{1.5}
  \setlength{\tabcolsep}{10pt}
  \begin{tabularx}{\textwidth}{|>{\bfseries\color{white}\columncolor{boaGreen}}l|X|}
    \hline
    \multicolumn{2}{|>{\bfseries\color{white}\columncolor{boaGreen}}c|}{Cas d'utilisation : Administration du système (CU-06)} \\ \hline
    Acteurs & Administrateur / Superviseur \\ \hline
    Pré-conditions & 
    L'administrateur s'est connecté et s'est authentifié sur le module Admin (AUTHENTIFICATION ADMIN). \\ \hline
    Scénario nominal & 
    1. L'administrateur accède au tableau de bord d'administration.\newline
    2. Le système affiche les indicateurs clés et l'état de l'automatisation.\newline
    3. L'administrateur choisit l'une des actions suivantes :\newline
    \quad a. Consulter le Tableau de Bord pour suivre les métriques.\newline
    \quad b. Consulter ou Supprimer des Historiques de discussion et de transactions.\newline
    \quad c. Gérer les clients BOA (Consulter / Ajouter / Supprimer un profil client).\newline
    4. Le système applique l'action et rafraîchit les vues. \\ \hline
    Scénarios alternatifs & 
    \textbf{A3. Échec de la suppression de client/historique :}\newline
    Si l'action de suppression échoue en base de données (ex: contrainte d'intégrité), le système affiche un message d'erreur et conserve les données intactes. \\ \hline
    Post-conditions & Les modifications administratives sont enregistrées en base de données et auditées. \\ \hline
  \end{tabularx}
  \caption{Description du cas d'utilisation - Administration}
  \label{tab:cu_admin}
\end{table}
